\section{The canonical comparison on each original affine open}

\begin{lemma}\label{lem:fc-module-flat-base-change-affine-open}
\lean{EnvelopeProperty.SchemeModuleFlatBaseChangeAffineOpen.isIso_comparison}
\leanok
\uses{scheme:existing-affine-module-pullback-sections,lem:fc-affine-cartesian-module-base-change,lem:fc-scalar-module-open-restriction,lem:fc-module-pullback-open-restriction,lem:fc-scheme-module-flat-base-change-map}
On the actual scalar base change of $p:X\to\operatorname{Spec}R$,
for any original quasicoherent $M$ and any original affine open $U$,
the canonical whole section comparison
$S\otimes_R\Gamma(U,M)\to\Gamma(g^{-1}U,g^*M)$ is an isomorphism.
\end{lemma}
\begin{proof}
The square obtained by restricting the original projection to $U$ and
$g^{-1}U$ is cartesian. Paste it with the actual original scalar base
change square. Its upper source and target are affine: the target is
$U$ and the source is its base change along the affine spectrum map.
The original quasicoherent module restricts to a quasicoherent module
on $U$. Apply the affine cartesian module comparison to this square.

Identify the original open-section modules with the entire global
section modules of the two original restrictions, retaining their
literal scalar actions. Original module pullback commutes with original
open restriction by uniqueness of the actual composite adjunctions.
The equality of the two composite adjunction units identifies the
affine comparison with the original canonical comparison on $U$:
apply the equality on $U$, transport the actual open-restriction unit
sections to the top opens, and use native presheaf naturality. On each
pure tensor both sides are the actual coefficient times the same
original pullback-unit section. Tensor extensionality gives equality
of the entire maps. The affine equivalence and the derived original
restriction isomorphisms therefore prove that the original canonical
component is an isomorphism. No compatibility square or section
comparison is supplied as an additional hypothesis.
\end{proof}
